\documentclass{amsart}

\usepackage{amsmath,amssymb,amsfonts,amscd}
\usepackage[all,cmtip]{xy}
\usepackage{enumerate}
\usepackage{ amssymb, latexsym, amsmath}
%\setcounter{MaxMatrixCols}{30}%
\usepackage{fullpage}
\usepackage{url}
\usepackage{hyperref}
\usepackage{ marvosym }
\usepackage{quiver}



\numberwithin{equation}{section}

%\usepackage{cjw-latex}

%\providecommand{\U}[1]{\protect\rule{.1in}{.1in}}

\theoremstyle{plain}
\newtheorem{theorem}{Theorem}[section]
\newtheorem*{thm}{Theorem}
\newtheorem{proposition}[theorem]{Proposition}
\newtheorem{lemma}[theorem]{Lemma}
\newtheorem{conjecture}[theorem]{Conjecture}
\newtheorem{corollary}[theorem]{Corollary}
\newtheorem{algorithm}[theorem]{Algorithm}
\newtheorem{axiom}[theorem]{Axiom}
\newtheorem{criterion}[theorem]{Criterion}
%\newtheorem{conjecture}[theorem]{Conjecture}
\newtheorem{wildconjecture}[theorem]{Wild Conjecture}


\theoremstyle{definition}
\newtheorem{definition}[theorem]{Definition}
\newtheorem*{dfn}{Definition}
\newtheorem{condition}[theorem]{Condition}
\newtheorem{example}[theorem]{Example}
\newtheorem{exercise}[theorem]{Exercise}
\newtheorem{notation}[theorem]{Notation}
\newtheorem{question}[theorem]{Question}
\newtheorem{problem}[theorem]{Problem}
\newtheorem{solution}[theorem]{Solution}




\theoremstyle{remark}
\newtheorem{remark}[theorem]{Remark}
\newtheorem{remarks}[theorem]{Remarks}
\newtheorem{summary}[theorem]{Summary}
\newtheorem{observation}[theorem]{Observation}
\newtheorem{conclusion}[theorem]{Conclusion}
\newtheorem{acknowledgement}[theorem]{Acknowledgement}
\newtheorem{case}[theorem]{Case}
\newtheorem{claim}[theorem]{Claim}



\makeatletter
\newcommand{\rmnum}[1]{\romannumeral #1}
\newcommand{\Rmnum}[1]{\expandafter\@slowromancap\romannumeral #1@}
\makeatother
\newcommand{\Aut}{\operatorname{Aut}}
\newcommand{\Ext}{\operatorname{Ext}}
\newcommand{\Tor}{\operatorname{Tor}}
\newcommand{\Z}{\mathbb{Z}}
\newcommand{\osum}{\oplus}
\newcommand{\aff}{\operatorname{Aff}}

\newcommand{\Id}{\operatorname{Id}}
\newcommand{\concat}{(\Id \to \Omega \Sigma)}
\newcommand{\eval}{(\Sigma\Omega \to \Id)}
\newcommand{\Q}{\mathbb{Q}}
\newcommand{\F}{\mathbb{F}}
\newcommand{\N}{\mathbb{N}}
\newcommand{\g}{\mathfrak{g}}
\newcommand{\n}{\mathfrak{n}}
\newcommand{\h}{\mathfrak{h}}
\newcommand{\p}{\mathfrak{p}}
\newcommand{\q}{\mathfrak{q}}
\newcommand{\m}{\mathfrak{m}}
\newcommand{\e}{\mathfrak{e}}
\newcommand{\f}{\mathfrak{f}}
\newcommand{\V}{\mathscr{V}}
\newcommand{\D}{\mathscr{D}}
\renewcommand{\S}{\mathscr{S}}
\newcommand{\Sch}{\mathscr{C}}
\newcommand{\Y}{\mathscr{Y}}
\renewcommand{\u}{\mathfrak{u}}
\newcommand{\set}[2]{ \left\{ {#1} \, \left| \, {#2} \right\}\right.}
\newcommand{\A}{\mathbb{A}}
\newcommand{\X}{\mathcal{X}}
\renewcommand{\Im}{\operatorname{Im}}
\newcommand{\T}{\mathscr{T}}
\newcommand{\HH}{\mathscr{H}}
\newcommand{\PGL}{\operatorname{PGL}}
\newcommand{\righthookarrow}{\hookrightarrow}
\newcommand{\sign}{\operatorname{sign}}
\renewcommand{\_}[2]{\underbrace{#1}_{#2}}
\renewcommand{\^}[2]{\overbrace{#1}_{#2}}
\newcommand{\curverightarrow}{\curvearrowright}

\newcommand{\z}{\mathfrak{z}}
\renewcommand{\sl}{\mathfrak{sl}}
\newcommand{\R}{\mathbb{R}}
\newcommand{\RP}{\mathbb{RP}}
\renewcommand{\P}{\mathbb{P}}
\newcommand{\C}{\mathbb{C}}
\renewcommand{\H}{\mathcal{H}}
\newcommand{\Ind}{\operatorname{Ind}}
\newcommand{\Tr}{\operatorname{Tr}}
\newcommand{\Orb}{\mathcal{O}}
\renewcommand{\k}{\mbox{\Fontauri k}}
\newcommand{\Ad}{\operatorname{Ad}}
\newcommand{\ad}{\operatorname{ad}}
\newcommand{\Hom}{\operatorname{Hom}}
\newcommand{\Ker}{\operatorname{Ker}}
\newcommand{\End}{\operatorname{End}}
\newcommand{\supp}{\operatorname{supp}}
\newcommand{\Stab}{\mbox{Stab}}
\newcommand{\Lie}{\operatorname{Lie}}
\newcommand{\Gal}{\operatorname{Gal}}
\newcommand{\GL}{\operatorname{GL}}
\newcommand{\SL}{\operatorname{SL}}
\newcommand{\Sp}{\operatorname{Sp}}
\newcommand{\Mp}{\operatorname{Mp}}
\renewcommand{\a}{\mathfrak{a}}
\newcommand{\Wh}{\operatorname{Wh}}
\newcommand{\Span}{\operatorname{Span}}
\newcommand{\WF}{\operatorname{WF}}
\newcommand{\AC}{\operatorname{AC}}
\newcommand{\Lin}{\operatorname{Lin}}
\newcommand{\Diff}{\operatorname{Diff}}
\newcommand{\Gen}{\operatorname{Gen\, Hom}}
\newcommand{\Her}{\operatorname{Her}}
\newcommand{\pivot}{&}
\newcommand{\sgn}{\operatorname{sgn}}
%\newcommand{C^{\infty}}{DS}
\newcommand{\vsp}{{\vspace{0.2in}}}



\title{Notes}

\author{Sudharshan K V}
\begin{document}
\maketitle

\section{Introduction}
Let $F$ be a local field. Roughly, the Local Langlands Correspondence (LLC) for $\GL_n/F$ is a \textbf{check} bijection between the $n$-dimensional (complex) representations of the Weil group $W_F$ of $F$ and the irreducible admissible (complex) representations of $\GL_n(F)$. The Weil group of a local field is a modification of its absolute Galois group. In this article, we will describe LLC for the local field $\R$. In fact, the archimedean case was the motivation for Langlands to consider such potential relations over non-archimedean local fields, leading to the formulation of LLC!\\

The representation theory of $W_\R$ is quite simple. On the other hand, the classification of irreducible admissible representations of a real reductive group upto infinitesimal equivalence was achieved by Langlands [cite], and is now called the Langlands classification. The LLC for $\GL_n/\R$ is evident from the Langlands classification for $\GL_n$, once the representations of $W_\R$ are identified. 

\section{The Weil group of $\mathbb R$}
For any local field $F$, one can define the Weil group $W_F$ axiomatically (see Tate). The Weil group for $\R$ is a non-split extension of $\Gal(\C/\R)$ by $\C^\times$. Concretely, we have $$W_\R = \C^\times \cup j\C^\times,$$ where $j^2 = -1, jzj^{-1} = \bar z$ for $z \in \C^\times$. These are the units of the quarternion algebra: the non-split extension is given by a non-trivial element of the Brauer group of $\R$ (which is $\Z/2\Z$), and this element corresponds to a non-split central simple algebra over $\R$ - the quarternions. We will however, work with the concrete description.

\section{Representations of $W_\R$}
In this section, we will describe the representations of $W_\R$. We will only consider (complex) representations of $W_\R$ where the group elements act semisimply.

\subsection{One dimensional representations}
The characters of $W_\R$ factor through its abelianization. One computes immediately the following:
\[jzj^{-1}z^{-1} = \bar z z^{-1},\] so the commutator subgroup is the unit circle $S^1$ in $\C^\times$. $\C^\times/S^1$ is the positive real numbers, and we can identify $jC^\times/S^1$ with the negative reals. Therefore, $W_\R^{\text{ab}} = \R^\times$. This is one of the axioms in [Tate].\\

The (continuous) characters of $\R>0$ just $x\mapsto x^s$ for $s\in \C$. The continuous characters of $\R^\times$ are thus parameterized by a complex variable $s$ and a character $\epsilon$ of $\{\pm 1\}$. We denote by $\chi_{s, \epsilon}$ the character which sends complex $z$ to $|z|^s$ and $j$ to $\epsilon(-1)$, where $|.|$ is the usual absolute value on $\C$. \\

Because of the identification $W_\R^\text{ab} \simeq \R^\times$, we have 

\begin{theorem}[LLC for $\GL_1/\R$] There is a bijection between irreducible (admissible) representations of $\GL_1(\R)$ and the one dimensional (semisimple) representations of the Weil group $W_\R$.
\end{theorem}

In general, any character of $\C^\times$ is given by a pair of complex numbers $\mu, \nu$, \[z \mapsto z^\mu \bar z^\nu,\]  where the only restriction is that $\mu - \nu$ is an integer. We denote this character by $\chi_{\mu,\nu}$. 

\subsection{Two dimensional representations}
Here, we identify the two dimensional irreducible representations of $W_\R$. As before, we assume that the elements of $W_\R$ act semisimply. \\

Let $\varphi$ be a two dimensional representation of $W_\R$. First, we restrict the representation $\varphi, V$ to $\C^\times$. The set of operators $\varphi(z), z \in \C^\times$ is a family of commuting semisimple operators, so we choose a basis $u,v$ of $V$ in which the matrix representation is diagonal and write \[\varphi =
  \begin{pmatrix}
    \chi_{\mu, \nu} & 0 \\
    0 & \chi_{\mu', \nu'}
  \end{pmatrix}.\]

Let $u' = \varphi(j) u$. For $z\in \C^\times$, we compute \[\varphi(z)(u') = \varphi(j\bar z j^{-1})(u') = \chi_{\mu, \nu}(\bar z)\varphi(j)(u)=z^\nu \bar z^\mu u'.\] So $u'$ is also a common eigenvector for the $\C^\times$ action. If the two characters $\chi_{\mu,\nu}$ and $\chi_{\mu',\nu'}$ are the same, then the representation is reducible - and decomposes as sum of eigenspaces of the semisimple operator $\varphi(j)$.
If $u'$ is a scalar multiple of $u$, then $\mu = \nu$, and similar reasoning for $v' = \varphi(j) v$ shows that $v' \in \C v$ and $\mu' = \nu'$. Again, this representation decomposes as a sum of two characters of $W_\R$.\\

Finally suppose that $u' \in \C v$ and that $\varphi$ is irreducible. In this case, we have $\mu = \nu'$ and $\mu' = \nu$ (and $\mu \neq \nu$). The representations restricting to $\chi_{\mu,\nu} \osum \chi_{\nu,\mu}$ and $\chi_{\nu,\mu}\osum \chi_{\mu,\nu}$ are isomorphic. Therefore, irreducible two dimensional representations of $W_\R$ are parameterized by a complex number $t = \frac12 \left(\mu +\nu\right)$ and a positive integer $l = \mu - \nu$. We denote this $\tau_{l,t}$

\subsection{Reducibility of representations}
Let $\varphi, V$ be a $n$-dimensional representation of $W_\R$. Consider the representation $V$ as a $\C^\times$ representation, and one has the decomposition into (one dimensional) eigenspaces \[V = \osum_{\mu,\nu} V_{\mu,\nu}.\] As before, one sees that the action of $\varphi(j)$ permutes the eigenspaces. If $\varphi(j)V_{\mu,\nu} = V_{\mu,\nu}$ then it is a one dimensional subrepresentation of $V$. Otherwise, $\varphi(j)V_{\mu,\nu} = V_{\nu,\mu}$, and as before, the representation $V_{\mu,\nu} \osum V_{\nu,\mu}$ is representation of $W_\R$.\\

\begin{theorem}
  The $n$ dimensional representations of $W_\R$ are completely reducible, and each irreducible constituent is either one or two dimensional.
\end{theorem}

\section{The Langlands Classification for $\GL_n(\R)$}
<<<<<<< HEAD
=======
In this section, we state Langlands classification for irreducible admissible representations of $\GL_n(\R)$ upto infinitesimal equivalence. As mentioned earlier, the representations of $\GL_1(\R)$ are parameterized by a complex variable $s$ and a character of $\{\pm 1\}$. We will first mention some special representations $\GL_2$ and then about how one builds representations of $\GL_n$ from these representations in rank $1$ and $2$.
>>>>>>> 11d626ee760b51f735940ef5fc1fa1524b8ca733

\subsection{Special Representations of $\GL_2$}

For $l\geq 2$ an integer, recall that $\SL_2(\R)$ has the discrete series representations $D_l^+$, and the limit of discrete series representations denoted $D_1^+$. Let $\SL_2^\pm(\R)$ denote the group of $2\times 2$ real matrices with determinant $\pm 1$. We consider the following (tempered) representations of $SL_2^\pm(\R)$: \[D_l = \Ind_{\SL_2(\R)}^{\SL_2^\pm(\R)}(D_l^+),\] for $l\geq 1$.\\
Next, consider the \emph{relative discrete series} representations $\sigma_{l,t}$ given by \[\sigma_{l,t} = D_l \otimes |\det|^t,\] where $l \geq 1$ and $t \in \C$. These representations will be the building blocks (in rank 2) for representations of $\GL_n(\R)$. (Relate this to the representation

\subsection{Langlands Classification for $\GL_n$}
We state Langlands classification of irreducible admissible representations of $\GL_n(\R)$ upto infinitesimal equivalence. There is a more general version for an arbitrary reductive group in place of $\GL_n$, but we will not state that here. See [getz hahn].

For any partition of $n = \sum_i^r n_i$ where $n_i = 1$ or $2$, and representations $\sigma_i$ of $\GL_{n_i}$ of the form mentioned above (i.e. characters of $\R^\times$ or the relative discrete series $\sigma_{l,t}$), let $M$ be the Levi subgroup $\prod_i \GL_{n_i}(\R)$ of $\GL_n(\R)$ and let $P$ be the standard parabolic of $\GL_n(\R)$ containing $M$. Inflate the representation $\sigma = (\sigma_1, \dots, \sigma_r)$ of $M$ to $P$, and consider the normalized induction \[I(\sigma_1,\dots,\sigma_r) = \Ind_P^{\GL_n(\R)} \left( (\sigma_1, \dots, \sigma_r) \otimes \frac12 \rho_P \right ),\] where $\rho_P$ denotes the sum of positive roots associated to the parabolic $P$. 

\begin{theorem}[Langlands]
\begin{enumerate}
\item Suppose $t_1/n_1 \geq \t_2/n_2 \geq \dots \geq t_r/n_r$, then $I(\sigma_1,\dots, \sigma_r)$ has a unique irreducible quotient denoted $J(\sigma_1,\dots,\sigma_r)$.
\item All irreducible admissible representations of $\GL_n(\R)$ are infinitesimally equivalent to some $J(\sigma_1, \dots, \sigma_r)$. 
\item $J(\sigma_1,\dots,\sigma_r)$ and $J(\sigma'_1,\dots,\sigma'_{r'})$ are infinitesimally equivalent if and only if $r = r'$ and the $\sigma_i'$ are a permutation of $\sigma_i$ (such that the condition in (1) holds).
\end{enumerate}
\end{theorem}

\begin{corollary}[LLC for $\GL_n/\R$]
There is a one-one correspondence between irreducible admissible representations of $\GL_n(\R)$ upto infinitesimal equivalence and the isomorphism classes of $n$-dimensional representations of $W_\R$ where the group acts semisimply.
\end{corollary}
\begin{proof}
Let $J = J(\sigma_1,\dots, \sigma_r)$ be an irreducible admissible representation of $\GL_n(\R)$. To each representation $\sigma_i$ of $\GL_{n_i}$, consider the associated $n_i$-dimensional Weil group representation: if $n_i = 1$, then consider the character $\tau_i: W_\R \to \R^\times \to \C^\times$, where the first map is the projection $W_\R \to W_\R^{\text{ab}}$ and the second map is $\sigma_i$. If $n_i = 2$, consider the representation $\tau_i = \tau_{l,t}$ described in Section \ref{sec:weil-rep}, where $l, t$ are the parameters associated to $\sigma_i$. Now associate to $J$ the $n$-dimensional representation $\sum_i \tau_i$ of $W_\R$.\\

This association is the desired bijection. The inverse map is clear: to any $n$-dimensional Weil representation $V$, consider its decomposition $V = \sum_{i=1}^r \tau_i$ into one and two dimensional pieces. For each $i$, consider the representation $\sigma_i$ of $\GL_1(\R)$ or $\GL_2(\R)$ with the same parameters as $\tau_i$ as above. Permute the factors so that the condition (1) of the theorem is satisfied, and consider $J(\sigma_1,\dots,\sigma_r)$.
\end{proof}

\end{document}
